% Generated draft: requires mathematical and source review.
\documentclass[11pt]{article}
\usepackage[margin=25mm]{geometry}
\usepackage{fontspec}
\setmainfont{Latin Modern Roman}
\usepackage{amsmath,amssymb,mathtools,unicode-math}
\setmathfont{Latin Modern Math}
\usepackage{enumitem,needspace,xurl,hyperref}
\setlength{\parindent}{0pt}
\setlength{\parskip}{5pt}
\setlength{\emergencystretch}{3em}
\newcommand{\N}{\mathbb{N}}
\newcommand{\Z}{\mathbb{Z}}
\newcommand{\Q}{\mathbb{Q}}
\newcommand{\R}{\mathbb{R}}
\newcommand{\C}{\mathbb{C}}
\newcommand{\F}{\mathbb{F}}
\newcommand{\A}{\mathbb{A}}
\newcommand{\PP}{\mathbb{P}}
\newcommand{\E}{\mathbb{E}}
\newcommand{\eps}{\varepsilon}
\newcommand{\dd}{\,\mathrm{d}}
\begin{document}
\section*{UnsolvedMath ID: 700710}
\textbf{Source problem code:} catalogue entry
\subsection{Definitions and mathematical statement}
Let \(G\) be a finite group. Use the commutator convention \([x,y]=x^{-1}y^{-1}xy\), and define the iterated commutators by
\[
[x,{}_0y]=x,\qquad [x,{}_{r+1}y]=[[x,{}_ry],y]\quad(r\geq 0).
\]
Here a law means an identity holding for every substitution \(x,y\in G\). Define
\[
d(G)=\min\bigl\{d\in\Z_{>0}:\exists r\in\Z_{\geq0}\ \forall x,y\in G,\ [x,{}_ry]=[x,{}_{r+d}y]\bigr\}.
\]
Thus \(r\) is uniform over all pairs \((x,y)\), and may depend on \(G\). For finite \(G\), this set of positive integers is nonempty. Let \(F(G)\) denote the Fitting subgroup, namely the largest nilpotent normal subgroup of \(G\). For a finite group \(H\), let \(\exp(H)\) be the least positive integer \(m\) such that \(h^m=1\) for every \(h\in H\).

\begin{enumerate}
\item[(a)] Set
\[
e(G)=\begin{cases}1,&d(G)\text{ is even},\\2,&d(G)\text{ is odd}.\end{cases}
\]
Is it true that every finite group \(G\) satisfies
\[
\exp\bigl(G/F(G)\bigr)\mid e(G)d(G)?
\]
\item[(b)] Is it true that every finite nonabelian simple group \(G\) satisfies
\[
\exp(G)\mid d(G)?
\]
\end{enumerate}
The supplied source states that part (a) is true for soluble groups, attributing this result to N. D. Gupta and H. Heineken. It also reports that the proposer checked part (b) for alternating groups \(\mathbb A_n\), groups \(\mathrm{PSL}(2,q)\), and a number of sporadic groups. These are supplied source qualifications, not independently verified assertions here. The source does not specify parameter ranges for those checks or identify the sporadic groups. It also does not explicitly define its commutator or zero-fold iteration notation; the conventions used above make that notation explicit and require review against the original source.
\subsection{Short English statement}
For a finite group, let \(d(G)\) be the smallest positive eventual period of the iterated commutators, with one starting index valid for every pair of elements. The questions ask whether the exponent modulo the Fitting subgroup divides \(d(G)\), allowing a factor of \(2\) when \(d(G)\) is odd, and whether the exponent divides \(d(G)\) itself for every finite nonabelian simple group. The supplied source reports the soluble case of (a) and checks of (b) for certain families.
\subsection{Sources}
\begin{itemize}
\item Local UnsolvedMath registry: permanent ID \(700710\), source code \texttt{catalogue entry}, title \emph{Kourovka Notebook Problem 11.17}. The supplied record labels the problem open; this label is not an independent verification of its current status.
\item R. Brandl, proposer, 1990. \emph{The Kourovka Notebook (Unsolved Problems in Group Theory)}, original problem 11.17, 11th issue (1990). The supplied background also identifies the 21st issue edition (2026), arXiv identifier \texttt{1401.0300v48}. The supplied collection description names E. I. Khukhro and V. D. Mazurov as editors.
\item N. D. Gupta and H. Heineken, \emph{Math. Z.} \textbf{95} (1967), 276--287. Supplied as the reference for part (a) for soluble groups; no article title is supplied.
\item The report of checks of part (b) for \(\mathbb A_n\), \(\mathrm{PSL}(2,q)\), and a number of sporadic groups is attributed to the proposer in the supplied statement; no separate reference for those checks is supplied.
\end{itemize}
\end{document}
